% Proposed isolated B.143. Full transfer proof; no private PDF/PNG.
\subsection*{Virtual nonsplitting and the exact fixed-source restriction}

\statement{B.143}{Proved}{Neither actual filling admits a virtual section, while its first-jet classes restrict to zero on the actual fixed source}
\label{stmt:B143}
Let \(\pi:\Gamma_0\to Q_0=T\rtimes S\) be B.142's
actual eight-family coarse filling. For every finite-index
subgroup \(H\leq Q_0\), the restriction
\(\pi^{-1}(H)\to H\) has no group-homomorphic right inverse.
For every open subgroup \(\mathcal H\leq\widehat Q_0\),
the restriction of the actual completed map has no continuous
group-homomorphic right inverse.

The same conclusions hold for B.139's actual one-family filling
\(p_a:\Gamma_0\to\Gamma_a\), over every finite-index subgroup
of \(\Gamma_a\) and every open subgroup of its full completion.
These are statements about whole-group virtual sections.

In contrast, write \(J_a=\Gamma_a^\delta\) for the actual
one-family fixed-source image, and \(\overline{J_a}\) for
its actual closure in \(\widehat\Gamma_a\).
B.142's integral first-jet class and every odd-prime class
pull back to zero on \(J_a\); every corresponding continuous
odd-prime class pulls back to zero on \(\overline{J_a}\).
The subgroups \(J_a\) and \(\overline{J_a}\) have infinite
index and are not open, respectively. No obstruction to a
fixed-source section follows from these particular classes.

\begin{proof}
\emph{B.143.1. The actual finite-coefficient input.}\par
B.142 constructs the actual one-color quotient
\[
 V=\mathbb Z[T]/I^2,\qquad V_p=V/pV,
\]
with its complete nontrivial \(Q_0\) action. Its actual extensions are
\begin{equation}\label{eq:B143-actual-extensions}
 \begin{gathered}
 1\longrightarrow V\longrightarrow E_V\longrightarrow Q_0
       \longrightarrow1,\\
 1\longrightarrow V_p\longrightarrow E_p\longrightarrow Q_0
       \longrightarrow1,\\
 1\longrightarrow V_p\longrightarrow\widehat E_p
       \longrightarrow\widehat Q_0\longrightarrow1.
 \end{gathered}
\end{equation}
The last kernel is the actual embedded finite subgroup,
and the extension admits a continuous set section.
For every odd p the classes
\[
 \omega_p\in H^2(Q_0,V_p),\qquad
 \widehat\omega_p\in H^2_{\rm cts}(\widehat Q_0,V_p)
\]
are nonzero: B.142's single eight-row certificate gives
\(0=2\) for every possible choice of retained and source lifts.
Its characteristic finite-kernel refinement supplies the
completed extension and its nonzero continuous class separately.
No comparison between discrete and continuous cohomology,
or completion of the infinite family kernel, is assumed.

\emph{B.143.2. Restriction with the full coefficient action.}\par
Here is the entire finite-index cohomological argument needed.
Let G be a discrete or profinite group, let \(H\leq G\)
have finite index n, open in the latter case, and let A be
an \(\mathbb F_p[G]\)-module, finite continuous in the
profinite case. If \(p\nmid n\), restriction
\[
                       H^2(G,A)\longrightarrow H^2(H,A)
 \tag{B.143.1}
\]
is injective, with continuous cohomology in the profinite case.
The action on A is not required to be trivial.

Use homogeneous cochains: a two-cochain \(c(g_0,g_1,g_2)\)
is G-equivariant, and its differential is the alternating
sum of omitted arguments. All cochains are continuous when
appropriate. Choose representatives t for the right cosets
\(H\backslash G\), choosing \(1\) for H. The map
\[
                 \rho:G\to H,\qquad \rho(ht)=h,
\]
is H-equivariant and is the identity on H. It is continuous
when H is open, since its finitely many cosets are clopen.

Suppose c is a cocycle whose restriction is \(\partial b\)
for an H-equivariant one-cochain \(b:H^2\to A\).
Put \(b_\rho(g_0,g_1)=b(\rho g_0,\rho g_1)\), and define
\[
 Kc(g_0,g_1)
   =c(\rho g_0,g_0,g_1)-c(\rho g_0,\rho g_1,g_1).
 \tag{B.143.2}
\]
Both cochains are H-equivariant. Expanding the cocycle
identity on the three four-point tuples gives
\[
       \partial Kc=c-\rho^*c,\qquad
       \partial b_\rho=\rho^*c.
 \tag{B.143.3}
\]
For example the first expansion replaces
\(-c(\rho g_0,g_0,g_2)+c(\rho g_0,g_0,g_1)\)
by \(c(g_0,g_1,g_2)-c(\rho g_0,g_1,g_2)\);
the tuples \((\rho g_0,\rho g_1,g_1,g_2)\) and
\((\rho g_0,\rho g_1,\rho g_2,g_2)\) cancel the remaining
terms to \(-c(\rho g_0,\rho g_1,\rho g_2)\).
Thus \(B=Kc+b_\rho\) satisfies \(\partial B=c\)
as an H-equivariant cochain on G.

Now choose representatives r for the left cosets \(G/H\)
and average with the actual coefficient action:
\[
 (\operatorname{Av}B)(g_0,g_1)
   =\sum_{r\in G/H}r\cdot
          B(r^{-1}g_0,r^{-1}g_1).
 \tag{B.143.4}
\]
Replacing r by rh has no effect, by H-equivariance of B.
Reindexing the finite sum proves G-equivariance. It is
continuous in the compact case, and
\[
 \begin{aligned}
 \partial(\operatorname{Av}B)(g_0,g_1,g_2)
   &=\sum_{r\in G/H}r\cdot
          c(r^{-1}g_0,r^{-1}g_1,r^{-1}g_2)\\
   &=n\,c(g_0,g_1,g_2).
 \end{aligned}
 \tag{B.143.5}
\]
The last equality is G-equivariance of c, not a
trivial-action assertion. Multiplication by n is invertible
in \(\mathbb F_p\), so \(n^{-1}\operatorname{Av}B\)
is a G-equivariant coboundary for c. This proves (B.143.1)
in both categories, including all signs and continuity.

\emph{B.143.3. No coarse virtual section.}\par
Let \(H\leq Q_0\) have arbitrary finite index n. Choose an
odd prime \(p\nmid n\). A section of the original filling
over H, composed with its actual quotient to \(E_p\),
would split the restricted extension in
\eqref{eq:B143-actual-extensions}. Its class would restrict
to zero on H. Part~2 would force \(\omega_p=0\),
contrary to part~1.

For an arbitrary open \(\mathcal H\leq\widehat Q_0\),
choose an odd p avoiding its finite index. A continuous
section of the original completed filling would similarly
split the restriction of \(\widehat E_p\), forcing
\(\widehat\omega_p=0\) by the continuous version of part~2.
This is again impossible. One uniform actual certificate
therefore handles all finite-index or open subgroups;
no primes, finite groups, or lift words are being enumerated.

\emph{B.143.4. The actual one-family virtual consequence.}\par
B.139 gives the actual maps
\[
 \Gamma_0\xrightarrow{p_a}\Gamma_a
       \xrightarrow{q_a}Q_0,\qquad q_ap_a=\pi,
 \qquad j_a:Q_0\to\Gamma_a,\quad q_aj_a=1,
 \tag{B.143.6}
\]
where \(j_a\) is the literal coarse section of its graph
semidirect product. Let \(H\leq\Gamma_a\) have finite index.
Then \(Q_H=j_a^{-1}(H)\) has finite index in \(Q_0\):
the map of left-coset sets
\[
        Q_0/Q_H\longrightarrow\Gamma_a/H,\qquad
                 qQ_H\longmapsto j_a(q)H
\]
is injective. A section of \(p_a^{-1}(H)\to H\)
composed with \(j_a|_{Q_H}\) would give a forbidden
coarse section over \(Q_H\).

The completed maps in (B.143.6) retain their identities;
in particular \(\widehat j_a\) is an actual embedding
with a continuous left inverse. For every open subgroup
\(\mathcal H\leq\widehat\Gamma_a\),
its inverse image under \(\widehat j_a\) is open.
A continuous one-family section over \(\mathcal H\)
would therefore contradict the completed coarse result.
No separability assumption on a separately completed
infinite subgroup is needed.

\emph{B.143.5. The classes vanish on the actual fixed-source restriction.}\par
B.139 proves that the actual discrete fixed source \(J_a\)
is generated by the graph norm \(A=aa^*\) and the actual
sources m,v. Its coarse projection sends A to one and
is the identity on \(S=F(m,v)\). Hence
\[
                 q_a(J_a)=S
\]
and \(q_a|_{J_a}\) factors through this actual S subgroup.
The extension \(E_V=(V\rtimes T)\rtimes S\) has a
literal homomorphic section on S, with the actual corrected
action on its first parenthesis. Write \(\operatorname{lift}_S\)
for that section. Its pullback over \(J_a\) has the explicit
homomorphic section
\[
              j\longmapsto
                 (j,\operatorname{lift}_S(q_a(j))).
 \tag{B.143.7}
\]
Thus the integral first-jet class, and all its finite
reductions, pull back to zero on \(J_a\).

For completion use the actual closed source image
\(\overline{J_a}\), not an asserted identification with
an abstract completion of \(J_a\). Continuity and density
give \(\widehat q_a(\overline{J_a})\subseteq\widehat S\);
equality holds because the closed source contains the
actual source generators and their embedded compact closure
\(\widehat S\). The latter embedding follows from the
literal source projection and section in
\(\widehat Q_0=\widehat T\rtimes\widehat S\).
B.142 also supplies the literal continuous source section
of \(\widehat E_p\). Therefore the actual pullback over
\(\overline{J_a}\) has the continuous section
\[
              z\longmapsto
                (z,\operatorname{lift}_{\widehat S}(\widehat q_a(z))).
 \tag{B.143.8}
\]
Every \(\widehat\omega_p\) pulls back to zero there.

Finally \(\Gamma_a/J_a\) maps onto \(Q_0/S\), whose
cosets are in bijection with the infinite retained group T.
Thus \(J_a\) has infinite index. The compact coset space
\(\widehat\Gamma_a/\overline{J_a}\) likewise maps onto
\(\widehat Q_0/\widehat S\), in bijection with the infinite
\(\widehat T\); hence \(\overline{J_a}\) is not open.
The finite-index transfer argument does not apply to this
restriction.

These are proved limitations of the obstruction, not merely
an absence of a proposed implication. A fixed-kernel argument
must retain further norm-direction information. Neither the
marked fixed-section condition nor the original fixed-source
equality is settled by the whole-group virtual nonsplitting
proved above.
\end{proof}

The complete prior actual transfer and restriction proofs are retained
with the genuine B.142 coefficient inputs:
\begin{quote}\footnotesize\raggedright
\path{work_A/actual_coarse_filling_no_virtual_section29.md},
\path{work_A/actual_one_family_no_virtual_section29.md},
\path{work_A/first_jet_fixed_source_zero_restriction30.md}.
\end{quote}
